\documentclass[11pt]{article}

\usepackage[margin=1.1in]{geometry}
\usepackage{amsmath,amssymb,amsthm,mathtools}
\usepackage{enumitem}
\usepackage{hyperref}
\usepackage[nameinlink,noabbrev]{cleveref}

\hypersetup{colorlinks=true,linkcolor=blue,citecolor=blue,urlcolor=blue}

\newtheorem{theorem}{Theorem}[section]
\newtheorem{lemma}[theorem]{Lemma}
\newtheorem{corollary}[theorem]{Corollary}
\theoremstyle{remark}
\newtheorem{remark}[theorem]{Remark}

\providecommand{\doi}[1]{\href{https://doi.org/#1}{doi:#1}}
\newcommand{\one}{\mathbf 1}

\title{Exact free norm bounds fail for permutation matrices at degree one}
\author{Yitzchak Shmalo\\
Einstein Institute of Mathematics, The Hebrew University of Jerusalem\\
Givat Ram, Jerusalem, Israel}
\date{}

\begin{document}
\maketitle

\begin{abstract}
Van Handel asked whether there are permutation matrices of arbitrarily large size $N$ with the following property,
for some $q(N)\to\infty$. On the space orthogonal to the constant vector, every noncommutative polynomial of degree
at most $q(N)$ has norm at most its free norm, that is, its norm in the left regular representation of the free
group. Such a
property would extend the Ramanujan property of regular graphs to all polynomials. We show that it fails already for
polynomials of degree one: for every $N\ge 2$, every $r\ge 2$ and every $r$-tuple of $N\times N$ permutation
matrices there is a self-adjoint polynomial of degree one, with scalar coefficients, whose norm on that space exceeds
its free norm by at least $1/(4(N-1)^4)$. The polynomial can be taken
from an explicit set of at most $2N(N+1)$ polynomials that depends only on $N$.
\end{abstract}

\noindent\textit{Mathematics Subject Classification (2020):} 46L54, 05C50, 15B51, 46L05.\\
\textit{Keywords:} strong convergence, permutation matrices, free group, Ramanujan graphs, operator norm.

\section{Introduction}\label{sec:intro}

Let $r\ge 2$, let $g_1,\dots,g_r$ be free generators of the free group $\mathbf F_r$, and let $u_i=\lambda(g_i)$,
where $\lambda$ is the left regular representation on $\ell^2(\mathbf F_r)$. For an $r$-tuple
$\mathbf U^N=(U_1^N,\dots,U_r^N)$ of $N\times N$ permutation matrices write $H_N=\one^\perp\subset\mathbb C^N$ for
the orthogonal complement of the constant vector; every permutation matrix and its adjoint preserve $H_N$. For a
noncommutative polynomial $P$ in the variables $x_1,\dots,x_r,x_1^*,\dots,x_r^*$ we write
$P(\mathbf U^N,\mathbf U^{N*})$ and $P(\mathbf u,\mathbf u^*)$ for the operators obtained by substituting
$U_i^N,U_i^{N*}$, respectively $u_i,u_i^*$, for $x_i,x_i^*$. We call $\|P(\mathbf u,\mathbf u^*)\|$ the \emph{free
norm} of $P$.
Bordenave and Collins \cite{BC} proved that for independent uniformly random permutation matrices
\[
  \lim_{N\to\infty}\bigl\|P(\mathbf U^N,\mathbf U^{N*})|_{H_N}\bigr\|=\|P(\mathbf u,\mathbf u^*)\|
\]
in probability, for every polynomial $P$ with matrix coefficients. This is the \emph{strong convergence} of
$\mathbf U^N|_{H_N}$ to $\mathbf u$. For the single polynomial $x_1+x_1^*+\dots+x_r+x_r^*$ the left side is the
norm on $H_N$ of the adjacency matrix of a $2r$-regular graph, and the free norm is $2\sqrt{2r-1}$, the norm of the
adjacency operator of the $2r$-regular tree. For this polynomial the theorem of Bordenave and Collins is Friedman's
theorem \cite{Fr} for this model of random $2r$-regular graphs: the norm of the adjacency matrix on $H_N$ tends to
$2\sqrt{2r-1}$ in probability.
A $d$-regular graph on $N$ vertices whose adjacency matrix has norm at most $2\sqrt{d-1}$ on $H_N$ is a Ramanujan
graph; such graphs exist for infinitely many degrees by the constructions of Lubotzky, Phillips and Sarnak
\cite{LPS} and of Margulis \cite{Ma}. A bipartite $d$-regular graph has the eigenvalue $-d$ on $H_N$; the bipartite
Ramanujan graphs of every degree of Marcus, Spielman and Srivastava \cite{MSS} satisfy the bound $2\sqrt{d-1}$ for
the remaining eigenvalues.

Van Handel \cite[Section~6.5]{vH} asks whether the strong convergence phenomenon has a Ramanujan form, namely
whether there are permutation matrices $\mathbf U^N$ for arbitrarily large $N$ such that
\begin{equation}\label{eq:vh}
  \bigl\|P(\mathbf U^N,\mathbf U^{N*})|_{H_N}\bigr\|\le\|P(\mathbf u,\mathbf u^*)\|
\end{equation}
for every polynomial $P$ of degree at most $q(N)$, where $q(N)\to\infty$ sufficiently slowly. He observes that
\eqref{eq:vh} cannot hold for all polynomials when $N$ is fixed, since the reduced group $C^*$-algebra of
$\mathbf F_r$ would then embed in the $N\times N$ matrices, and writes that he is not aware of an obstruction under a
degree bound. He also poses the weaker question in which $\mathbf U^N$ may depend on the polynomial $P$.

The obstruction for fixed $N$ uses all polynomials at once and gives no bound on the degree.
O'Donnell and Wu \cite[Theorem~10.13]{OW} fix a degree $k$, a bound $R$ on the coefficients and
$\varepsilon>0$. For every large $N$ they construct an explicit permutation tuple of a size between $N$ and
$N+o(N)$ that serves every self-adjoint polynomial $P$ of degree at most $k$ with coefficients bounded by $R$ at
once: the spectrum of $P(\mathbf U^N,\mathbf U^{N*})|_{H_N}$ lies within Hausdorff distance $\varepsilon$ of the
spectrum of $P(\mathbf u,\mathbf u^*)$. So all these polynomials satisfy \eqref{eq:vh} up to the additive error
$\varepsilon$ at one tuple. Alev and Rao \cite[Theorem~5.5]{AR} take a fixed polynomial $Q$ with nonnegative
coefficients and evaluate it at adjacency matrices of regular graphs of bounded degree, each conjugated by a
permutation matrix. For every choice of the permutations, the norm on $H_N$ is at least the norm of $Q$ at free
copies of the adjacency operators of the corresponding regular trees, minus an error that tends to zero as
$N\to\infty$.

The inequality \eqref{eq:vh} holds with equality for polynomials of degree zero. It fails already for polynomials
of degree one, for every $N\ge 2$ and every tuple. Put $n=N-1$ and write $\mathbf U=(U_1,\dots,U_r)$ for
$\mathbf U^N$. For indices $i\ne j$, a complex
number $z$ with $|z|=1$, a sign $\sigma\in\{1,-1\}$ and $t>0$ let
\begin{equation}\label{eq:P}
  P=\bar z\,x_i+z\,x_i^*+\sigma t\,(x_j+x_j^*),
\end{equation}
and call $\|P(\mathbf U,\mathbf U^*)|_{H_N}\|-\|P(\mathbf u,\mathbf u^*)\|$ the \emph{gap} of $P$ at $\mathbf U$.
The free norm of $P$ does not depend on $i$, $j$, $z$ or $\sigma$ (\Cref{lem:phase}), and
\begin{equation}\label{eq:freebound}
  \|P(\mathbf u,\mathbf u^*)\|\ \le\ 2+2t^{4/3}
\end{equation}
(\Cref{lem:free}). The exact value is given by the formula of Akemann and Ostrand \cite{AO} for weighted sums of
free generators, in the form stated by Bordenave and Lacoin \cite[(62)]{BL}; \Cref{lem:free} proves the upper
bound by a weighted Schur test. For any two $N\times N$ permutation matrices, one of them, $C$, has a unit eigenvector
$v\in H_N$, $Cv=zv$, at which the other, $D$, satisfies
\[
  |\operatorname{Re}\langle v,Dv\rangle|\ \ge\ \frac1n
\]
(\Cref{lem:finite}). With $U_i=C$, $U_j=D$ and $\sigma$ the sign of $\operatorname{Re}\langle v,Dv\rangle$,
\[
  \langle v,P(\mathbf U,\mathbf U^*)v\rangle=2+2t\,|\operatorname{Re}\langle v,Dv\rangle|\ \ge\ 2+\frac{2t}n ,
\]
and $2t/n>2t^{4/3}$ when $t<n^{-3}$. The gap is then positive for every $N\ge 2$. In \cite{AR} the norm on $H_N$ is
bounded below by the free norm minus an error term, for a fixed polynomial with nonnegative coefficients. Here it
is bounded below by the free norm plus $1/(4n^4)$, for a polynomial that depends on the tuple and whose
coefficients $\bar z$, $z$ and $\sigma t$ need not be nonnegative. In \cite{OW} one tuple serves all polynomials
of bounded degree and bounded coefficients up to the error $\varepsilon$; here no tuple serves all polynomials of
degree one with no error.

\begin{theorem}\label{thm:main}
Let $N\ge 2$, $r\ge 2$ and $n=N-1$, and let $\mathbf U=(U_1,\dots,U_r)$ be $N\times N$ permutation matrices.
There are indices $i\ne j$, a complex number $z$ with $|z|=1$ and a sign $\sigma\in\{1,-1\}$ such that the
polynomial \eqref{eq:P} with $t=1/(2n^3)$ satisfies
\[
  \bigl\|P(\mathbf U,\mathbf U^*)|_{H_N}\bigr\|\ \ge\ 2+\frac1{n^4}\ \ge\ \|P(\mathbf u,\mathbf u^*)\|+\frac1{4n^4}.
\]
In particular, for $r\ge 2$ and $N\ge 2$ no $r$-tuple of $N\times N$ permutation matrices satisfies \eqref{eq:vh}
for all polynomials of degree at most one.
\end{theorem}

The polynomial depends on $\mathbf U$, as the question allows. The number $z$ is an eigenvalue of a permutation
matrix, hence a root of unity of order at most $N$, and this gives a set that does not depend on $\mathbf U$.

\begin{corollary}\label{cor:family}
Let $N\ge 2$, $r\ge 2$ and $n=N-1$, and let $\mathcal P_N$ be the set of polynomials \eqref{eq:P} with $t=1/(2n^3)$,
$(i,j)\in\{(1,2),(2,1)\}$, $\sigma\in\{1,-1\}$ and $z^k=1$ for some $1\le k\le N$. Then $|\mathcal P_N|\le 2N(N+1)$,
and for every $r$-tuple $\mathbf U$ of $N\times N$ permutation matrices some $P\in\mathcal P_N$ satisfies
$\|P(\mathbf U,\mathbf U^*)|_{H_N}\|\ge\|P(\mathbf u,\mathbf u^*)\|+1/(4n^4)$.
\end{corollary}

\begin{remark}\label{rem:scope}
\begin{enumerate}[label=(\alph*)]
\item The lower bound $1/(4n^4)$ for the gap tends to zero, so \Cref{thm:main} is consistent with strong convergence and with the
constructions of \cite{OW}. For $t=1/(2n^3)$ the coefficients of \eqref{eq:P} have $\ell^1$ norm $2+2t\le 3$, so
after normalization the gap is at least $1/(12n^4)$.
\item For $r=1$ the inequality \eqref{eq:vh} holds for every Laurent polynomial $P$ and every unitary matrix $U$:
the spectrum of $U$ lies on the unit circle, so $\|P(U,U^*)\|=\max_{\zeta\in\operatorname{spec}U}|P(\zeta,\bar\zeta)|
\le\max_{|\zeta|=1}|P(\zeta,\bar\zeta)|=\|P(u,u^*)\|$, the last equality by the Fourier transform on $\ell^2(\mathbb Z)$.
So two generators are needed.
\item \Cref{lem:finite} uses that the matrices are permutation matrices. In \Cref{thm:main} the polynomial is
chosen after the tuple; the weaker question, in which the tuple is chosen after the polynomial, is not settled by it.
\end{enumerate}
\end{remark}


\section{The free side}\label{sec:free}

This section bounds the free norm of the polynomial \eqref{eq:P}. Let $a,b$ be free generators of $\mathbf F_2$
and, for $t>0$,
\[
  T_t=\lambda(a)+\lambda(a)^*+t\bigl(\lambda(b)+\lambda(b)^*\bigr),
\]
which for $r=2$ is $P(\mathbf u,\mathbf u^*)$ with $i=1$, $j=2$, $z=1$ and $\sigma=1$. The exact value of $\|T_t\|$
is given by the formula of Akemann and Ostrand \cite{AO}, in the form of \cite[(62)]{BL}; only the upper bound
below is used, and it has the following direct proof.

\begin{lemma}\label{lem:free}
For all $s,t>0$,
\[
  \|T_t\|\ \le\ \sqrt{s^2+4}+\sqrt{s^2+4t^2}-s\ \le\ 2+\frac{s^2}4+\frac{2t^2}s .
\]
\end{lemma}

\begin{proof}
Put $\rho_a=(\sqrt{s^2+4}-s)/2$ and $\rho_b=(\sqrt{s^2+4t^2}-s)/(2t)$. Then $\rho_a^{-1}=(\sqrt{s^2+4}+s)/2$ and
$\rho_b^{-1}=(\sqrt{s^2+4t^2}+s)/(2t)$, so that
\begin{equation}\label{eq:rho}
  \rho_a^{-1}-\rho_a=s,\qquad t\,(\rho_b^{-1}-\rho_b)=s .
\end{equation}
Join each $w\in\mathbf F_2$ to $gw$ for $g\in\{a^{\pm1},b^{\pm1}\}$; this graph is a $4$-regular tree. Give the
edge $\{w,gw\}$ the weight $c_{w,gw}=1$ if $g=a^{\pm1}$ and $c_{w,gw}=t$ if $g=b^{\pm1}$, and call it an $a$-edge or
a $b$-edge accordingly. Then $T_t$ is the weighted adjacency operator of this tree. Let $h(w)>0$ be the product, over the letters of the reduced word $w$, of $\rho_a$ for each letter
$a^{\pm1}$ and $\rho_b$ for each letter $b^{\pm1}$; so $h(e)=1$. Write $w'\sim w$ when $w'$ and $w$ are joined, and $(T_th)(w)=\sum_{w'\sim w}c_{w,w'}h(w')$. At the
root, $(T_th)(e)/h(e)=2\rho_a+2t\rho_b$. At a vertex $w\ne e$ whose first letter is $a^{\pm1}$, the neighbours are the
parent, one further vertex across an $a$-edge and two across $b$-edges, so by \eqref{eq:rho}
\[
  \frac{(T_th)(w)}{h(w)}=\rho_a^{-1}+\rho_a+2t\rho_b=2\rho_a+2t\rho_b+s ,
\]
and the same value is obtained when the first letter is $b^{\pm1}$. Hence $(T_th)(w)\le\mu\,h(w)$ for every $w$, with
$\mu=2\rho_a+2t\rho_b+s=\sqrt{s^2+4}+\sqrt{s^2+4t^2}-s$.

For finitely supported $f$, summing $2|f(w)f(w')|\le|f(w)|^2h(w')/h(w)+|f(w')|^2h(w)/h(w')$ over the edges
$\{w,w'\}$ with weight $c_{w,w'}$ gives
\[
  |\langle f,T_tf\rangle|\le\sum_w|f(w)|^2\,\frac{(T_th)(w)}{h(w)}\le\mu\|f\|^2 .
\]
Since $T_t$ is bounded and self-adjoint and finitely supported vectors are dense, $\|T_t\|\le\mu$. The second
inequality of the lemma follows from $\sqrt{s^2+4}\le 2+s^2/4$ and $\sqrt{s^2+4t^2}\le s+2t^2/s$.
\end{proof}

With $s=(4t^2)^{1/3}$ the last bound of \Cref{lem:free} is $2+3\cdot 4^{-1/3}t^{4/3}\le 2+2t^{4/3}$, which with
\Cref{lem:phase} gives \eqref{eq:freebound}. For $t=1/(2n^3)$ this choice is $s=1/n^2$, the value used in the proof of \Cref{thm:main}.

\begin{lemma}\label{lem:phase}
Let $r\ge 2$, $i\ne j$, $|z|=1$, $\sigma\in\{1,-1\}$ and $t>0$. For the polynomial $P$ of \eqref{eq:P},
$\|P(\mathbf u,\mathbf u^*)\|=\|T_t\|$.
\end{lemma}

\begin{proof}
There is a homomorphism $\chi$ from $\mathbf F_r$ to the unit circle with $\chi(g_i)=z$, $\chi(g_j)=\sigma$ and
$\chi(g_k)=1$ otherwise. The diagonal unitary $(M\xi)(w)=\chi(w)\xi(w)$ on $\ell^2(\mathbf F_r)$ satisfies
$(M\lambda(g)M^*\xi)(w)=\chi(w)\chi(g^{-1}w)^{-1}\xi(g^{-1}w)=\chi(g)(\lambda(g)\xi)(w)$, that is,
$M\lambda(g)M^*=\chi(g)\lambda(g)$, and hence $M\lambda(g)^*M^*=\chi(g)^{-1}\lambda(g)^*$. So $M$ conjugates $\bar z u_i+zu_i^*+\sigma t(u_j+u_j^*)$ to $u_i+u_i^*+t(u_j+u_j^*)$. The subgroup
$K$ generated by $g_i,g_j$ is free on these two generators, and $\ell^2(\mathbf F_r)$ is the orthogonal direct sum of
the spaces $\ell^2(Kx)$ over the right cosets $Kx$ of $K$. Each of them is invariant under $\lambda(g_i)$ and
$\lambda(g_j)$, and on it these two operators are unitarily equivalent to the left regular representation of
$\mathbf F_2$, with $g_i,g_j$ in place of $a,b$. Hence the norm equals $\|T_t\|$.
\end{proof}

\section{The finite side}\label{sec:finite}

This section bounds the norm of $P(\mathbf U,\mathbf U^*)$ on $H_N$ from below and proves \Cref{thm:main}.
A permutation matrix $U$ permutes the basis vectors $\delta_x$ of $\mathbb C^N$; a fixed point or a cycle of $U$
is a fixed point or a cycle of this permutation.

\begin{lemma}\label{lem:finite}
Let $N\ge 2$ and let $A,B$ be $N\times N$ permutation matrices. There are an ordering $(C,D)$ of $(A,B)$, a
unit vector $v\in H_N$ and a complex number $z$ with $|z|=1$ such that
\[
  Cv=zv\qquad\text{and}\qquad|\operatorname{Re}\langle v,Dv\rangle|\ge\frac1n .
\]
\end{lemma}

\begin{proof}
Suppose first that one of the two matrices, call it $C$, fixes a point $x$, and let $D$ be the other. Put
\[
  v=\Bigl(1-\frac1N\Bigr)^{-1/2}\Bigl(\delta_x-\frac1N\one\Bigr).
\]
Then $v\in H_N$, $\|v\|=1$ and $Cv=v$. Since $D\one=D^*\one=\one$ and $\langle\delta_x,D\delta_x\rangle$ is $1$ if
$D$ fixes $x$ and $0$ otherwise,
\[
  \langle v,Dv\rangle=\frac{\langle\delta_x,D\delta_x\rangle-1/N}{1-1/N},
\]
which equals $1$ or $-1/n$.

Suppose now that neither matrix has a fixed point. The restriction $A|_{H_N}$ is unitary; let
$v_1,\dots,v_n$ be an orthonormal basis of $H_N$ consisting of its eigenvectors. Then
\[
  \sum_{k=1}^n\operatorname{Re}\langle v_k,Bv_k\rangle=\operatorname{Re}\operatorname{Tr}\bigl(B|_{H_N}\bigr)
  =\operatorname{Tr}B-1=-1,
\]
since $B$ has no fixed point. Some summand is at most $-1/n$, and that $v_k$ works with $(C,D)=(A,B)$.
\end{proof}

\begin{proof}[Proof of \Cref{thm:main}]
Apply \Cref{lem:finite} to $U_1,U_2$; this gives $(C,D)=(U_i,U_j)$ with $\{i,j\}=\{1,2\}$, the vector $v$ and $z$.
Choose $\sigma$ with $\sigma\operatorname{Re}\langle v,Dv\rangle>0$. Since $C$ is unitary, $C^*v=\bar zv$, so
\[
  \langle v,P(\mathbf U,\mathbf U^*)v\rangle=\bar z z+z\bar z+2\sigma t\operatorname{Re}\langle v,Dv\rangle
  \ \ge\ 2+\frac{2t}n=2+\frac1{n^4}.
\]
The operator $P(\mathbf U,\mathbf U^*)$ is self-adjoint and preserves $H_N$, so its norm there is at least this
Rayleigh quotient. By \Cref{lem:phase} and \Cref{lem:free} with $s=1/n^2$,
\[
  \|P(\mathbf u,\mathbf u^*)\|=\|T_t\|\le 2+\frac1{4n^4}+\frac{2t^2}{s}=2+\frac1{4n^4}+\frac1{2n^4}
  =2+\frac3{4n^4},
\]
and the gap is at least $1/(4n^4)$.
\end{proof}

\begin{proof}[Proof of \Cref{cor:family}]
Every eigenvalue of a permutation matrix is a $k$-th root of unity, where $k$ is the length of a cycle, so $k\le N$.
Hence the polynomial constructed in the proof of \Cref{thm:main} lies in $\mathcal P_N$. There are at most
$\sum_{k=1}^Nk=N(N+1)/2$ roots of unity of order at most $N$, two choices of $(i,j)$ and two signs.
\end{proof}

\begin{remark}\label{rem:example}
For $N=3$ let $C$ be the matrix of the transposition of $1$ and $2$ (fixing $0$) and $D$ the matrix of the
cycle $0\to1\to2\to0$. Take
$\mathbf U=(C,D)$, $t=1/16$ and $P=x_1+x_1^*-t(x_2+x_2^*)$, the polynomial \eqref{eq:P} with $z=1$ and $\sigma=-1$.
On $H_3$ we have $D+D^*=-I$, so $P(\mathbf U,\mathbf U^*)$ acts there as $2C+1/16$: the vector $(2,-1,-1)$ is an
eigenvector with eigenvalue $33/16$ and $(0,1,-1)$ is one with eigenvalue $-31/16$, and
$\|P(\mathbf U,\mathbf U^*)|_{H_3}\|=33/16$. \Cref{lem:phase} and \Cref{lem:free} with $s=1/4$ give
$\|P(\mathbf u,\mathbf u^*)\|\le 131/64$, so the gap is at least $1/64$. For the same two matrices the
adjacency matrix $C+C^*+D+D^*$ has norm $3<2\sqrt3$ on $H_3$, so the corresponding $4$-regular graph is Ramanujan.
\end{remark}

\section*{Acknowledgments and funding}

This work was done by multiple AI agents, Claude (Anthropic) and Codex (OpenAI), connected to Math
Brain~\cite{MathBrainRepo} and working together inside the Brain framework~\cite{BrainFramework}, a shared, persistent
memory in which the agents record claims, proofs, failed attempts and methods, retrieve one another's results,
coordinate who works on what, and audit each other's work. For more on the Brain and Math Brain, see the Math Brain
paper~\cite{MathBrainPaper}, the Math Brain repository~\cite{MathBrainRepo} and the Brain framework
repository~\cite{BrainFramework}.

The research presented in this paper was supported by the European Research Council (ERC) under the European
Union's Horizon Europe research and innovation programme (grant agreement No.\ 101041711). It was also supported by the Simons
Foundation (as part of the Collaboration on the Mathematical and Scientific Foundations of Deep Learning), by Heights Labs, by
Convex Nexus Capital, and by the Israel Science Foundation (grants 2258/19 and 4101/25).

\section*{Declaration of competing interest}

The author holds equity in Heights Labs and Convex Nexus Capital and has a management role at Convex Nexus Capital.

\section*{Declaration of generative AI and AI-assisted technologies}

The arguments of \Cref{sec:free,sec:finite} were developed with AI systems (OpenAI Codex and Anthropic Claude agents
working under the author's direction), which also drafted the text and wrote the computer checks. The question is van
Handel's \cite{vH}; the free norm estimate is an instance of the formula of Akemann and Ostrand \cite{AO}, and the
setting is the strong convergence theorem of Bordenave and Collins \cite{BC}. The author directed the work and is
responsible for the content of the paper and for decisions concerning its dissemination.

\section*{Data and code availability}

A Python script that checks \Cref{lem:finite} and the inequality
$\|P(\mathbf U,\mathbf U^*)|_{H_N}\|\ge 2+1/n^4$ of \Cref{thm:main}, for the polynomial of the proof and for some
member of $\mathcal P_N$, for every ordered pair of permutation matrices of size $N\le 5$, and the example of
\Cref{rem:example} in exact arithmetic, is available at
\url{https://github.com/yspennstate/permutation-degree-one-obstruction}.

\begin{thebibliography}{99}

\bibitem{AO}
C.~A.~Akemann and P.~A.~Ostrand,
\emph{Computing norms in group $C^*$-algebras},
Amer.\ J.\ Math.\ \textbf{98} (1976), 1015--1047.
\doi{10.2307/2374039}

\bibitem{AR}
V.~L.~Alev and S.~Rao,
\emph{Rotating walks, derandomized squares, and permutation conjugated walks},
arXiv:2611.13693v2, 2026.

\bibitem{BL}
C.~Bordenave and H.~Lacoin,
\emph{Cutoff at the entropic time for random walks on covered expander graphs},
J.\ Inst.\ Math.\ Jussieu \textbf{21} (2022), 1571--1616.
\doi{10.1017/S1474748020000663}; arXiv:1812.06769.

\bibitem{BC}
C.~Bordenave and B.~Collins,
\emph{Eigenvalues of random lifts and polynomials of random permutation matrices},
Ann.\ of Math.\ (2) \textbf{190} (2019), 811--875.
\doi{10.4007/annals.2019.190.3.3}

\bibitem{Fr}
J.~Friedman,
\emph{A proof of Alon's second eigenvalue conjecture and related problems},
Mem.\ Amer.\ Math.\ Soc.\ \textbf{195} (2008), no.~910.
\doi{10.1090/memo/0910}

\bibitem{LPS}
A.~Lubotzky, R.~Phillips and P.~Sarnak,
\emph{Ramanujan graphs},
Combinatorica \textbf{8} (1988), 261--277.
\doi{10.1007/BF02126799}

\bibitem{Ma}
G.~A.~Margulis,
\emph{Explicit group-theoretic constructions of combinatorial schemes and their applications in the construction of
expanders and concentrators},
Problemy Peredachi Informatsii \textbf{24} (1988), no.~1, 51--60.

\bibitem{MSS}
A.~W.~Marcus, D.~A.~Spielman and N.~Srivastava,
\emph{Interlacing families I: Bipartite Ramanujan graphs of all degrees},
Ann.\ of Math.\ (2) \textbf{182} (2015), 307--325.
\doi{10.4007/annals.2015.182.1.7}

\bibitem{OW}
R.~O'Donnell and X.~Wu,
\emph{Explicit near-fully X-Ramanujan graphs},
in: 2020 IEEE 61st Annual Symposium on Foundations of Computer Science (FOCS), 2020, 1045--1056;
full version arXiv:2009.02595.
\doi{10.1109/FOCS46700.2020.00101}; arXiv:2009.02595.

\bibitem{vH}
R.~van Handel,
\emph{The strong convergence phenomenon},
arXiv:2507.00346v2, 2026.

\bibitem{MathBrainRepo}
Y.~Shmalo, \emph{Math Brain: public snapshot, code and tests}, 2026.
\url{https://github.com/yspennstate/math-brain-public}

\bibitem{MathBrainPaper}
Y.~Shmalo, \emph{Math Brain: shared, auditable memory for multi-agent mathematical research}, October 2026.
\url{https://github.com/yspennstate/math-brain-public/tree/main/papers/math-brain-2026-10-24}

\bibitem{BrainFramework}
Y.~Shmalo, \emph{Brain Framework}, public software repository, 2026.
\url{https://github.com/yspennstate/brain-framework-public}

\end{thebibliography}

\end{document}
